
% =====================================================================
\section{T2 and T3: the residual is a conserved region}
\label{sec:residual}
% =====================================================================

We now identify the quantity that survives every re-encoding of a contact graph.
The parts (vertices) may be relabelled, merged, re-described; what cannot be
removed is the boundary content between them---the \emph{cut-residual}. We prove
it is bounded below (T2), realised by no single vertex so that it is a region and
never a point (T2), and invariant under every weighted isomorphism (T3), hence
the unique conserved object from which the relabellable surface derives.

\subsection{The cut-residual}

\begin{definition}[Cut-residual of a partition]
\label{def:residual}
For a partition $\mathcal{P}=\{U_1,\dots,U_k\}$ of $\V$ into nonempty parts, its
\emph{cut-residual} is the total weight of inter-part boundary,
\[
\Res(\mathcal{P})\;=\;\tfrac12\sum_{i=1}^{k}\wt(\cut(U_i))
\;=\!\!\sum_{\substack{i<j}}\!\!\wt(\cut(U_i,U_j)),
\]
the sum over distinct part-pairs of the weight of the edges crossing between
them. The \emph{residual of $\Gph$} is the minimum over nontrivial partitions,
\[
\Res(\Gph)\;=\;\min_{\mathcal{P}\,:\,k\ge 2}\Res(\mathcal{P}),
\]
attained by finiteness of $\V$.
\end{definition}

\begin{remark}[The residual is the ``interstice'']
\label{rem:interstice}
$\Res(\mathcal{P})$ is the content of the boundaries \emph{between} the parts:
not what any part contains, but what lies in the gaps that separate them. As the
parts are packed (the partition refined or coarsened), this interstitial content
changes shape; T3 shows its total, suitably taken, does not change under mere
relabelling. The residual is the graph's ``negative space.''
\end{remark}

\subsection{T2: the residual is bounded below and is a region}

\begin{theorem}[Residual positivity and non-locality]
\label{thm:residual-region}
Let $\Gph$ be a contact graph with floor $\floorc>0$. Then:
\begin{enumerate}[label=\textup{(\roman*)},leftmargin=2.2em]
\item \textbf{Bounded below.} Every nontrivial partition has
$\Res(\mathcal{P})\ge\floorc>0$; hence $\Res(\Gph)\ge\floorc>0$.
\item \textbf{Realised by no single vertex.} There is no vertex $v$ with
$\Res(\Gph)=\wt(\cut(\{v\}))$ forced; the residual is a property of the
partition (a region of the graph), not of any one vertex. Equivalently, the
minimiser of \Cref{def:residual} is in general not a singleton, and the residual
cannot be localised to a point.
\end{enumerate}
\end{theorem}

\begin{proof}
(i) A nontrivial partition has $k\ge2$ parts; since $\Gph$ is connected, at
least one part-pair $(U_i,U_j)$ has $\cut(U_i,U_j)\neq\varnothing$, contributing
weight $\ge\floorc$ by \Cref{thm:floor}. Hence $\Res(\mathcal{P})\ge\floorc>0$,
and taking the minimum, $\Res(\Gph)\ge\floorc$.

(ii) Suppose, for contradiction, the residual were always realised by isolating a
single vertex, i.e. $\Res(\Gph)=\min_v\wt(\cut(\{v\}))$ for every contact graph.
Consider any graph in which the minimum-weight nontrivial cut separates a part
$U$ with $\abs{U}\ge 2$ from its complement at total weight strictly below
$\min_v\wt(\cut(\{v\}))$---for instance two dense clusters joined by a single
floor-weight edge, where the cheapest cut splits the two clusters
($\wt=\floorc$) but every singleton has several incident edges
($\wt(\cut(\{v\}))\ge 2\floorc$). Then the residual is realised by the
bipartition, not by any singleton, contradicting the supposition. Hence the
residual is in general a property of a multi-vertex partition: it is a region,
not a point.
\end{proof}

\begin{corollary}[No point carries the residual; the sorites for parts]
\label{cor:point-forbidden}
The residual of a contact graph is never carried by a lone vertex in general:
``what separates the parts'' is borne by a region of positive boundary content,
not by any single part. Consequently there is no sharp vertex at which a coarse
aggregate becomes a distinct part---the transition is spread across a band of
boundary content of width $\ge\floorc$, the graph analogue of the sorites
\citep{black1937,williamson1994}.
\end{corollary}

\begin{proof}
The first claim is \Cref{thm:residual-region}(ii). For the second: growing a
part one vertex at a time changes its boundary cost $b(U)$ by a bounded amount
per step (the weights of the edges whose crossing status flips, each $\ge
\floorc$ and finite in number), so $b(U)$ cannot jump across any threshold in a
single step from a value bounded away from it; the part becomes distinguishable
across a range of stages, of width controlled by $\floorc$, not at a point.
\end{proof}

\subsection{T3: the residual is invariant under re-encoding}

We now show the residual is a graph invariant: it depends on the graph up to
weighted isomorphism, not on any labelling. Combined with the fact that the
vertices (the surface data) are freely relabellable, this makes the residual the
conserved object beneath the surface.

\begin{theorem}[Conservation of the residual]
\label{thm:residual-invariant}
The residual $\Res(\Gph)$ is a weighted-graph invariant: if
$\phi\colon\Gph\to\Gph'$ is a weighted isomorphism then
$\Res(\Gph)=\Res(\Gph')$. In particular $\Res$ is invariant under every
automorphism of $\Gph$, hence under every relabelling of its vertices.
\end{theorem}

\begin{proof}
Let $\phi\colon\Gph\to\Gph'$ be a weighted isomorphism. For any partition
$\mathcal{P}=\{U_1,\dots,U_k\}$ of $\V$, the image
$\phi(\mathcal{P})=\{\phi(U_1),\dots,\phi(U_k)\}$ is a partition of $\V'$, and
$\phi$ being edge- and weight-preserving gives, for each pair,
$\wt'(\cut'(\phi(U_i),\phi(U_j)))=\wt(\cut(U_i,U_j))$. Summing,
$\Res(\phi(\mathcal{P}))=\Res(\mathcal{P})$. The correspondence
$\mathcal{P}\leftrightarrow\phi(\mathcal{P})$ is a bijection between nontrivial
partitions of $\V$ and of $\V'$ (with inverse $\phi^{-1}$), so the minima agree:
$\Res(\Gph')=\min_{\mathcal{P}'}\Res(\mathcal{P}')=
\min_{\mathcal{P}}\Res(\phi(\mathcal{P}))=\min_{\mathcal{P}}\Res(\mathcal{P})=
\Res(\Gph)$.
\end{proof}

\begin{theorem}[The residual is the conserved invariant beneath the surface]
\label{thm:residual-conserved}
Let $\mathcal{R}(\Gph)$ denote the orbit of $\Gph$ under all weighted
relabellings of its vertices \textup{(}its isomorphism class\textup{)}. Then:
\begin{enumerate}[label=\textup{(\roman*)},leftmargin=2.2em]
\item every vertex label is variable across $\mathcal{R}(\Gph)$---no labelling is
fixed; and
\item $\Res$ is constant on $\mathcal{R}(\Gph)$ and strictly positive.
\end{enumerate}
Hence, among data attached to a contact graph, the residual is an invariant of
the isomorphism class while the vertex labels are not: it is conserved under
exactly the operation---re-encoding---that moves everything else.
\end{theorem}

\begin{proof}
(i) For any vertex $v$ and any other vertex $v'$ of the same degree-and-weight
type, there is a relabelling carrying $v$ to a different label; more simply, the
identity of a vertex is a label, and labels are by definition permuted across the
isomorphism class. So no label is fixed across $\mathcal{R}(\Gph)$. (ii) is
\Cref{thm:residual-invariant} (constancy) together with
\Cref{thm:residual-region}(i) (positivity). The two together say: the
relabellable surface (vertex identities) carries no invariant, while the residual
is invariant and positive.
\end{proof}

\begin{remark}[Knowing is repacking, not draining]
\label{rem:repacking}
\Cref{thm:residual-conserved} has a reading we record but do not rely on. If one
identifies the vertices (the parts one has individuated) with ``what is known''
and the residual (the conserved interstice) with ``what remains true regardless
of how one has carved things up,'' then acquiring or re-describing knowledge is
the relabelling and repartitioning of vertices---which by
\Cref{thm:residual-invariant} leaves the residual unchanged in total. One never
drains the residual by re-encoding; one only re-shapes which boundaries one
holds. The conserved residual is, in this reading, the invariant ``truth'' that
no re-encoding removes, bounded below by the floor and---by
\Cref{thm:residual-region}(ii)---never a point. The reading is offered as
orientation only; the theorems above are purely graph-theoretic.
\end{remark}

\begin{remark}[A definite-yet-uncloseable quantity]
\label{rem:cat-parity}
A residual can be perfectly well-defined and yet impossible to fix by any single
completed partition. Consider a quantity such as the parity of $\abs{\V}$ when
the vertex set is itself in flux---vertices entering and leaving faster than any
partition can be committed. The parity is definite at every instant yet never
stabilises long enough to be read off a fixed graph; a finite procedure
correctly leaves it in the residual rather than paying unboundedly to chase it.
This illustrates \Cref{thm:residual-region}(ii): the residual is what a finite
process rationally leaves uncommitted, the ``sufficient unknown,'' and forcing it
to a point would cost without return. (Offered as illustration; no theorem
depends on it.)
\end{remark}
